\chapter{Emotional Lability}

\section*{Definition}
Emotional lability (also called pseudobulbar affect or pathological laughter and crying) refers to involuntary, excessive, or uncontrollable episodes of laughing, crying, or other emotional expressions that are out of proportion or incongruent with the triggering stimulus.

\section*{Pathophysiology}
Emotional lability arises from disruption of the cerebellar--brainstem--prefrontal circuits that normally modulate volitional emotional expression. Lesions in the corticopontine-cerebellar pathways (common in neurological disease) impair this regulation, disinhibiting the brainstem laughing/crying centers (nucleus ambiguus, facial nucleus) from voluntary control.

\section*{Causes}
\begin{itemize}
  \item \textbf{Neurological:} Pseudobulbar affect (most common causes: ALS, multiple sclerosis, stroke, traumatic brain injury, Alzheimer disease, Parkinson disease, brain tumors)
  \item \textbf{Psychiatric:} Bipolar disorder (rapid cycling), borderline personality disorder, post-traumatic stress disorder, histrionic personality traits
  \item \textbf{Medication-induced:} Antidepressants (especially during titration), levodopa, dopaminergic agonists, stimulants, withdrawal from SSRIs
  \item \textbf{Hormonal:} Premenstrual dysphoric disorder, perimenopause, postpartum period
  \item \textbf{Infectious:} Syphilis (general paresis), encephalitis
\end{itemize}

\section*{When to See a Doctor}
See your doctor if:
\begin{itemize}
  \item Episodes are frequent, intense, and incongruent to the situation
  \item There is a known neurological condition with new onset of lability
  \item Emotional outbursts are causing social, occupational, or relationship impairment
\end{itemize}

\section*{Related Symptoms}
\begin{itemize}
  \item Pathological laughter/crying, mood swings, irritability, disinhibition, euphoria, cognitive impairment
\end{itemize}
\textbf{Important:} Emotional lability in neurological disease (pseudobulbar affect) is distinct from mood disorders and can be misdiagnosed as depression.

\section*{Self-Care / Home Management}
\begin{itemize}
  \item Recognize triggers and use distraction or cognitive reframing to reduce episode intensity
  \item Practice deep breathing during episodes to regain composure
  \item Communicate with close contacts about the involuntary nature of the symptoms
  \item Avoid alcohol and sleep deprivation, which can worsen lability
\end{itemize}

\section*{How Your Doctor Will Evaluate You}
\begin{itemize}
  \item Differentiate pseudobulbar affect from bipolar disorder, borderline personality, and depression
  \item Screening tools: Center for Neurologic Study--Lability Scale (CNS-LS), Pathological Laughter and Crying Scale (PLACS)
  \item Neurological examination to identify underlying brain injury or neurodegeneration
  \item MRI brain if no known neurological diagnosis
\end{itemize}

\section*{Treatment Options}
\begin{itemize}
  \item First-line for pseudobulbar affect: dextromethorphan/quinidine (Nuedexta), a combination NMDA antagonist and sigma-1 agonist
  \item Off-label: SSRIs (citalopram, fluoxetine) or tricyclic antidepressants (amitriptyline) for milder cases
  \textit{For psychiatric causes: mood stabilizers (lithium, lamotrigine) for bipolar disorder; dialectical behavior therapy for personality disorders}
  \item Speech and language therapy for compensatory communication strategies
\end{itemize}

\clearpage
